% Figure: the third law and the vanishing of heat capacity as T -> 0. The Debye
% model gives c_v proportional to T^3 at low temperature for a solid; an electronic
% contribution adds a term linear in T in a metal. Either way the heat capacity is
% driven to zero as T -> 0, which is required by the third law because the entropy
% integral S(T) = integral_0^T c_p/T dT' must converge. We plot the Debye T^3 law
% rising from zero and the high-temperature Dulong-Petit plateau it approaches.
\begin{figure}[htbp]
\centering
\begin{tikzpicture}
\begin{axis}[
  width=11cm, height=7cm,
  xlabel={reduced temperature $T/\theta_D$},
  ylabel={molar heat capacity $c_v/(3R)$},
  xmin=0, xmax=1.6, ymin=0, ymax=1.1,
  axis lines=left,
  legend style={font=\scriptsize, at={(0.97,0.05)}, anchor=south east, draw=enggray!40},
  tick label style={font=\scriptsize},
  label style={font=\small},
  samples=160,
]
% Debye-like interpolation: low-T T^3, high-T Dulong-Petit plateau at 1.
% Use a smooth surrogate c/(3R) = x^3/(x^3 + a) tuned so it -> 1 and -> ~x^3.
% with a small constant a it reproduces both limits qualitatively.
\addplot[engblue, very thick, domain=0.001:1.6]
  {(x^3)/(x^3 + 0.10)};
\addlegendentry{$c_v/3R$ (Debye-type solid)}
% the pure cubic low-T asymptote
\addplot[engred, thick, dashed, domain=0.001:0.6]
  {10*x^3};
\addlegendentry{low-$T$ asymptote $c_v\propto T^3$}
% the Dulong-Petit plateau
\addplot[enggray, dotted, very thick, domain=0:1.6] {1};
\addlegendentry{Dulong-Petit plateau $3R$}
\node[engblue, font=\scriptsize, anchor=west] at (axis cs:1.0,0.85)
  {classical plateau};
\node[engred, font=\scriptsize, anchor=west] at (axis cs:0.30,0.18)
  {$c_v\to 0$ as $T\to 0$};
\end{axis}
\end{tikzpicture}
\caption[Heat capacity vanishing as the third law requires]{The molar heat
capacity of a solid, normalised to its classical Dulong-Petit value $3R$, plotted
against temperature scaled by the Debye temperature $\theta_D$. At high
temperature the capacity reaches the Dulong-Petit plateau, where every vibrational
mode is classically excited and equipartition holds. As the temperature falls the
modes freeze out one by one and the capacity drops; in the Debye model it falls as
$T^3$ near absolute zero, the red asymptote. The third law forces this vanishing:
the absolute entropy $S(T) = \int_0^T (c_p/T')\,dT'$ can converge to a finite
value at $T=0$ only if $c_p$ goes to zero at least as fast as $T$, so a material
with a non-vanishing heat capacity at absolute zero would carry infinite entropy,
which the count $S = k_B \ln W$ forbids for a system with a unique ground state.
The freezing-out of heat capacity at low temperature is the macroscopic shadow of
the discreteness of the quantum energy levels.}
\label{fig:vol04ch04:third-law-cp}
\end{figure}
